\documentclass[11pt,a4paper]{article}
\usepackage[margin=2.5cm]{geometry}
\usepackage{amsmath,amsfonts,amssymb}
\usepackage{graphicx}
\usepackage{authblk}
\usepackage{hyperref}
\usepackage{caption}
\usepackage{subcaption}
\usepackage{xspace}
\usepackage{float}
\usepackage{booktabs}
\usepackage{siunitx}
\newcommand{\hilumi}{HL-LHC}
\newcommand{\amcnlo}{\textsc{MadGraph5\_amc@nlo}}
\newcommand{\mgfive}{MG5aMC}
\newcommand{\ttbar}{\ensuremath{\mathrm{t}\bar{\mathrm{t}}}}
\newcommand{\gluglu}{\ensuremath{\mathrm{gg}}}
\newcommand{\ggtt}{\ensuremath{\gluglu\rightarrow\ttbar}\xspace}
\newcommand{\ggttj}{\ensuremath{\gluglu\rightarrow\ttbar\mathrm{j}}\xspace}

\title{\textbf{Accelerating Monte Carlo Event Generation using Hardware Acceleration Platforms}} 

\author[1]{J. Fern\'andez Men\'endez}
\author[2]{L. Fiorini}
\author[1]{S. Folgueras}
\author[2]{F. Herv\'as \'Alvarez}
\author[2]{H. Guti\'errez Arance}
\author[1]{P. Leguina L\'opez}
\author[2]{A. Oyanguren}
\author[2]{A. Valero}
\author[3]{C. Vico Villalba}

\affil[1]{Universidad de Oviedo - ICTEA}
\affil[2]{Instituto de F\'isica Corpuscular (IFIC) - U. Valencia}
\affil[3]{Rice University}


\date{}

\begin{document}
\maketitle

\begin{abstract}
We present a heterogeneous hardware approach to accelerate Monte Carlo event generation using the MadGraph5\_aMC@NLO generator. Targeting the matrix element and color matrix computation steps, we implement key components using High-Level Synthesis (HLS), VHDL, and AI Engine programming on an AMD Versal platform. Two benchmark processes---$e^+e^- \to \mu^+\mu^-$ and $pp \to t\bar{t} +$ jets---demonstrate performance gains and energy efficiency compared to conventional CPU and GPU implementations. The results showcase the scalability of FPGA acceleration for high-throughput and sustainable simulation workloads in high-energy physics.
\end{abstract}

\section{Introduction} % ---- From Santi (removed by Carlos)
% Physics event generators are a critical component of simulations in high-energy physics, used to produce synthetic collision events for both theory studies and experimental analyses. As the field moves into the High Luminisity LHC (\hilumi) era, the demand for simulated events is skyrocketing~\cite{HSFPhysicsEventGeneratorWG:2020gxw}, while traditional CPU-based computing improvements have slowed. Projections indicate that flat computing budgets and the end of Moore's law speedups could leave future experiments limited by Monte Carlo (MC) statistics if generator performance is not dramatically improved. 

% Need to add references to LSZ formula and so on probably
Due to the intrinsically probabilistic nature of fundamental particle interactions and the composite structure of the proton, accurately modelling the dynamics of proton collisions at the Large Hadron Collider (LHC) constitutes one of the major challenges in high-energy physics (HEP). Physics event generators encode the necessary tools to describe these collisions in a programmatic way, enabling the use of computational resources to make quantitative predictions for a given theoretical model.

Each collision at the LHC is characterized by an initial and a final state. The particles in the initial state originate from the colliding protons, whereas the final-state particles are produced as a result of the interaction. It can be shown that the transition probability (\(\mathcal{P}\)) for an initial state \(i\), consisting of two incoming partons \(q_1\) and \(q_2\), to evolve into a final state \(\mathcal{X}\) is proportional to the squared matrix element \(\mathcal{M}\), as shown in Equation~\ref{eq:hard_scatter_bracket}.

\begin{equation}
    \mathcal{P}(i \rightarrow \mathcal{X}) = |\mathcal{M}_{i \rightarrow \mathcal{X}}|^2
    \label{eq:hard_scatter_bracket}
\end{equation}

To compute the probability of a given process \(\mathcal{X}\) occurring at the LHC, one must integrate the squared matrix element over the full phase space that connects the initial and final states. These calculations can only be performed up to a fixed order in a perturbation theory that expands over different powers of the so-called coupling constant, that quantifies the strength of the interaction between particles. To do so, physics event generators rely on Monte Carlo (MC) integration techniques, which approximate the integral by sampling random points that populate the accessible phase space. All modern MC generators provide calculations at leading order (LO), while some are capable of higher-order predictions such as next-to-LO (NLO) and next-to-NLO (NNLO). These higher order corrections are not negligible in most of the cases, and with the ever increasingly size of the LHC dataset, experiments at the LHC have become sensitive to these corrections. However, increasing the perturbative order leads to a rapid, non-linear growth in computational complexity, placing severe demands on computing resources.

% ---- From Santi (removed by Carlos)
% In complex LHC processes, the parton-level event generation stage can consume a large fraction of the overall simulation time – for example, the calculation of matrix elements (the squared quantum amplitudes for a given process) can account for on the order of 30–40\% of the total event generation CPU time in multi-jet processes


% This is because for each simulated event, the generator must evaluate a complicated physics matrix element function at a random point in phase space, integrating over many Feynman diagram contributions. Generating millions of events thus entails repeating these calculations millions of times. In principle, this task is embarrassingly parallel: the same matrix element evaluation is performed independently for many different phase-space points. This parallelism has motivated efforts to accelerate event generators using modern hardware architectures. Notably, GPUs have been applied to MadGraph5\_aMC@NLO~\cite{Alwall:2014hca}, yielding substantial speedups in matrix element computation (often by one to two orders of magnitude for complex processes) by exploiting their massive parallel throughput~\cite{Valassi:2021ljk,Valassi:2025xfn,Hageboeck:2023blb,Valassi:2023yud}.

As a result, as the volume of data recorded by the LHC experiments continues to grow, so do the requirements for accurate MC simulations. To address this growing computational burden, developers of MC generators have devoted significant effort to parallelizing matrix element integration using vectorized CPU and GPU infrastructures.

While GPUs and CPUs have thus demonstrated the value of heterogeneous computing for event generation~\cite{Carrazza:2021gpx}, the use of Field-Programmable Gate Arrays (FPGAs) remains relatively unexplored in this domain~\cite{Barbone:2023mul}. FPGAs allow custom hardware architectures tailored to the specific computation, offering the potential for fine-grained parallelism and pipeline concurrency that can drastically speed up repetitive arithmetic-heavy tasks. Importantly, FPGAs often operate at lower clock frequencies but with much lower overhead, translating to excellent energy efficiency. For certain workloads, FPGAs have been shown to be significantly more energy-efficient than GPUs – in some cases over an order of magnitude in performance-per-watt. This makes FPGA acceleration an attractive approach to sustainable high-throughput computing for HEP, where power and cooling constraints are growing concerns. However, historically FPGAs have been challenging to program for complex algorithms like those in MC generators, which involve hundreds of operations and complex control flow generated from Feynman rules. Recent advances in high-level synthesis (HLS) tools and heterogeneous devices aim to lower these barriers. In particular, the AMD Versal Adaptive Compute Acceleration Platform (ACAP) is a new FPGA-based system-on-chip that integrates not only reconfigurable logic but also an array of specialized AI Engine cores and that promises to deliver exceptional compute density with low latency and power usage. The Versal platform thus provides a versatile heterogeneous playground to map both custom logic and parallel DSP workloads, and it has been demonstrated to deliver ``high throughput, low latency, and power efficient'' processing for suitable applications. 
% Insert the introduction content here

%The increasing complexity of HEP experiments, the growing volume of data, the improvement on the precision of the experimental measurements, and the progress on the theoretical calculations have placed significant demands on Monte Carlo (MC) simulations. The physics program of the Large Hadron Collider (LHC) experiments, require generating billions of events with an ever increasing need for more precise theoretical calculations, demanding significant processing time and resources. Traditional CPU-based computing platforms, while effective for many tasks, are facing limitations in terms of performance and scalability. 

%To address these limitations, accelerating computing platforms, such as Graphics Processing Units (GPUs), Tensor Processing Units (TPUs), have emerged as promising solutions. These platforms offer significantly higher computational power and parallel processing capabilities compared to CPUs, making them well-suited for the demanding tasks involved in MC simulations.  Field Programmable Gate Arrays (FPGAs) are becoming increasingly popular, and thus ofgreatly increasing relevance in the context of High-Performance Computing (HPC). Recent advances in High-Level Synthesis (HLS) toolchains for customised hardware implementations greatly reduce the engineering effort needed to program FPGAs. As a consequence, the use of FPGA accelerators is quickly spreading from the HPC context to other disciplines that require processing large volumes of data or ultra-low latency response.  By leveraging the parallel architecture of these accelerators, it is possible to achieve significant performance improvements and reduce simulation times and cost.  

%Monte Carlo simulations are widely utilised in the context of HEP, combined with their computational complexity they account for a large fraction of the total computational resources needed in HEP. In recent years, the increasing complexity of such simulations require significant computational time in order to accurately calculate the scattering amplitudes and matrix elements. In the last years the parallelisation of such calculations have been explored by frameworks such as \textit{MadFlow}\cite{Barbone:2023mul}, achieving great performance improvement while running on GPU. Some works has also been done demonstrating the possibility of achieving even better performance by using FPGAs. Barbone et al.\cite{Carrazza:2021gpx} demonstrate the use of FPGA acceleration by achieving an 100x speedup compared to multi-core CPU implementations. 

\section{MadGraph5\_aMC@NLO as a benchmark for hardware acceleration}


% ---- From Santi (Removed by Carlos)
% MadGraph5\_aMC@NLO~\cite{Alwall:2014hca,Frederix:2018nkq} (MG5aMC) is a state-of-the-art framework widely used in high-energy physics for the automatic generation of matrix elements and event simulation at leading (LO) and next-to-leading order (NLO) in perturbative quantum field theory. It provides a modular, extensible, and community-supported platform that is actively being adapted to meet the computational challenges posed by the HL-LHC and future colliders.

% Although experimental uncertainties today are larger than those of simulations, at the High Luminosity LHC experimental precision is expected to be above the theoretical one for events generated below NLO precisions. As forecasted hardware resources will not meet CPU requirements for these simulation needs, speeding up NLO event generation is a necessity.

% The integration of GPU acceleration has demonstrated substantial performance improvements. For complex processes such as $gg\rightarrow t\bar{t}ggg$, the achieved speed-up on the ME calculation utilization of NVIDIA V100 GPUs has resulted in ME computation ranging from 100x to 1000x compared to single-threaded Fortran implementations. This acceleration translates to an overall workflow speed-up of approximately 60x to 100x, considering the entire event generation pipeline. Additionally, the use of vector instructions on CPUs has led to ME calculation speed-ups of up to 16x. When combined with multi-core processing, the overall event generation workflow experiences a performance enhancement of approximately 6x to 10x.

% NLO computations introduce additional complexities, including real emission processes and loop corrections, which inherently involve more intricate algorithms and data dependencies compared to LO calculations. To address these challenges, the MG5aMC development team is implementing a multi-event processing framework that facilitates vectorization and parallel execution. This approach aims to exploit the capabilities of modern hardware architectures, such as GPUs, to accelerate the evaluation of NLO amplitudes. Preliminary efforts have demonstrated the feasibility of this strategy, with ongoing work focusing on optimizing the computational kernels and integrating them seamlessly into the existing MG5aMC infrastructure~\cite{Wettersten:2025hrb}. 

% On this work we will focus on the generation of the simple $gg\rightarrow t\bar{t}g$ process. It is true that in this process the non-ME part is dominating the CPU computation, but the ME-part is simple enough that allows for a simple case-study that we show in this manuscript. 


Among the most widely used frameworks for MC event generation in high-energy physics is the \amcnlo~(\mgfive) software. The team has already made significant strides in acceleration methods, particularly focusing on GPU and vector CPU support. The release of the first production version of the CUDACPP plugin (v1.00.00) ~\cite{Valassi:2025xfn} enables the acceleration of LO processes using GPUs and vector CPUs and involves extensive re-engineering to support data-parallel execution, facilitating significant speed-ups in ME computations, which traditionally represent the primary computational bottleneck in event generation workflows. 

The integration of GPU acceleration has demonstrated substantial performance improvements. For complex processes, the achieved speed-up on the ME calculation utilization of NVIDIA V100 GPUs has resulted in ME computation ranging from 100x to 1000x improvements compared to single-threaded Fortran implementations. This acceleration translates to an overall workflow speed-up of approximately 60x to 100x, considering the entire event generation pipeline. Additionally, the use of vector instructions on CPUs has led to ME calculation speed-ups of up to 16x. When combined with multi-core processing, the overall event generation workflow experiences a performance enhancement of approximately 6x to 10x.

The \mgfive~framework provides several features that make it particularly suitable for use as a benchmark when testing the performance of matrix element integration on hardware-based accelerators such as FPGAs. One such feature is the high degree of flexibility available to users to configure physics processes, model parameters, and integration settings. Once these parameters are fixed, the \mgfive~solution provides an automatic way of performing the matrix element integration by taking care of producing the necessary code to run the actual calculations. As a result, evaluating the limits of resource usage, throughput, numerical precision, and latency of an FPGA-based solution can be practically achieved by porting the matrix element code generated by \mgfive~to a hardware description language (HDL) of choice. A key advantage of this approach is that it removes the need for in-depth knowledge of the specific physics calculations being performed, since \mgfive~already provides validated and optimized implementations of the required matrix elements. This allows the focus to be placed directly on the hardware performance and architectural aspects of the accelerator, rather than on the underlying theoretical formulation of the problem.

As this work represents a first attempt at implementing complex physical calculations on a hardware accelerator, the focus is placed on exploring the limitations of FPGAs in one of the simplest scenarios that can be realized in hardware: the production of a \ttbar~final state through the interaction of two gluons from the incoming colliding protons, namely \(\ggtt\). The accuracy of the calculation is restricted to the LO perturbative expansion, since the complexity of NLO calculations can be challenging even for GPU-based solutions. To increase the complexity of the problem and provide a more demanding benchmark for the FPGA, we nevertheless include the first real-emission corrections, which constitute a subset of the full set of contributions required to achieve true NLO accuracy. To this end, we instruct \mgfive~to generate the code for the \ggttj~process, where the $\mathrm{j}$ denotes an additional emission from either the initial-state (\gluglu) or final-state (\ttbar) particles.

Within \mgfive, the complete integration of the matrix element is factorized into different steps. The most demanding step for processes such as \ggttj correspondes to the computation of the color algebra [...] % need to fill the gap here.






\section{Methodology} 
In this paper we present a first feasibility study of accelerating MC generation using FPGAs. We will stablish a comparison for a given physics process on the computing speed and power consumption between CPU, GPU and FGPA platforms. 

The following hardware platforms are considered: 
\begin{itemize}
    \item An \textbf{AMD Versal AI Core series device}, which combines traditional FPGA fabric with an array of AI Engine cores on a single chip. The Versal’s programmable logic is utilized for custom datapaths, while the AI Engines (tightly-coupled VLIW vector processors) are exploited for parallel numeric computations. The design was developed and tested on a Versal development board (VCK190).
\end{itemize}

The performance of these devices will be compared to GPU/CPU implementations that already exist within the MG5aMC project. 

\subsection{Target Hardware Platform}
We are using the Xilinx Versal AI Core Series VCK190 Evaluation Kit, which represents the inaugural evaluation platform for the Versal AI Core family, around the XCVC1902 Adaptive Compute Acceleration Platform (ACAP) device. 

The device features a dual-core ARM Cortex-A72 application processor capable of reaching 1.7 GHz, providing high single-threaded performance for system control and compute-intensive tasks. Complementing this is a dual-core ARM Cortex-R5F real-time processor with functional safety certification options, enabling deterministic, low-latency control for demanding real-time applications. 

The VC1902 device contains 400 AI Engine cores operating at 1 GHz or higher, delivering a peak compute throughput of 133 INT8 TOPs (tera-operations per second). These AI Engines are specialized processors designed specifically for machine learning inference and advanced signal processing, organized in a tightly-coupled array with deterministic, low-latency data movement and configurable memory hierarchies that avoid traditional cache miss penalties. Each AI Engine tile comprises: 
\begin{itemize}
    \item A 512-bit SIMD (Single Instruction Multiple Data) vector processing unit supporting both fixed-point and floating-point arithmetic,
    \item A 32-bit scalar RISC processor for control flow,
    \item 32 KB of local data memory with access to 128 KB of addressable shared memory through neighbouring tile connections,
    \item Dual 256-bit load units and single store units with individual address generators
    \item 16 KB of dedicated program memory-
\end{itemize}

The 1,968 DSP58 Engines provide complementary computational capacity optimized for high-precision floating-point (FP32/FP16) and fixed-point (INT8/INT16/INT24) arithmetic, delivering 2.8 TFLOP/s peak floating-point throughput and 11.8 TOP/s INT8 performance. These DSP engines are distributed across the programmable logic in dedicated DSP columns, enabling simultaneous deployment of DSP and general-purpose logic kernels.

The VCK190 provides 1,968,000 system logic cells, 899,840 LUTs (lookup tables), and extensive embedded memory: 158 MB UltraRAM and Block RAM distributed across dedicated memory columns; and 
4 MB Accelerator RAM shareable across compute engines for ultra-low-latency tensor storage. 

Additionally the device incorporates high-speed transceiver arrays and comprehensive memory capabilities. 

\subsection{Event Implementation}
The typical workflow of any matrix element generator, and MG5aMC in particular implies the following steps. The event's phase space is sampled using a pseudo-random number generator to generate the particles' momenta. Then, for each momenta, the matrix element, which is the probability of the event for the given process, is calculated. Traditionally, these steps occur sequentially, but can, however be calculated in parallel for many events at a time~\cite{Wettersten:2023ekm}. 

The Matrix Element computation is, by far, the most computationally intensive bottleneck. As it has been demonstrated~\cite{Valassi:2025xfn}, it is a highly-parallelizable process and thus resulting in significant speeds-up of the computation. In this work we will explore the feasiblity of achieving such acceleration using ACAP devices. 

\subsubsection{Matrix Element computation}

\subsubsection{Color algebra calculation}

\subsection{Test Processes and Workflows}

\subsection{Benchmark Metrics and Tools}


\section{Results}

\subsection{Performance Comparison}
\subsection{Power and Energy Efficiency}
\subsection{Scalability and Resource Utilization}


\section{Conclusion}

\section*{Acknowledgments}
This work was supported by [Funding Source], under grant/project number [XXX]. The authors acknowledge the contributions of [collaborators or institutions].

\bibliographystyle{unsrt}
\bibliography{references}

\end{document}


